\begin{table}[htbp]
\centering\small
\caption{What error correction buys, per controlled perturbation.}
\label{tab:qr_ec}
\begin{tabular}{lrrrrr}
\toprule
Transformation & L & M & Q & H & L$-$H (pp) \\
\midrule
\texttt{motion} & 100.0 & 100.0 & 100.0 & 100.0 & $+0.0$ \\
\texttt{blur} & 100.0 & 99.7 & 99.3 & 98.4 & $+1.5$ \\
\texttt{saltpepper} & 90.3 & 84.9 & 80.2 & 75.9 & $+14.4$ \\
\texttt{invert} & 71.6 & 68.9 & 69.0 & 67.5 & $+4.1$ \\
\texttt{perspective} & 66.5 & 62.6 & 59.2 & 56.6 & $+10.0$ \\
\texttt{logo} & 81.6 & 52.1 & 33.7 & 5.7 & $+75.9$ \\
\texttt{rotate} & 14.5 & 10.6 & 8.1 & 6.9 & $+7.6$ \\
\texttt{contrast} & 1.8 & 1.3 & 0.8 & 0.4 & $+1.4$ \\
\texttt{clean} & 0.2 & 0.3 & 0.3 & 0.4 & $-0.2$ \\
\bottomrule
\end{tabular}
\\[4pt]
\begin{minipage}{0.86\linewidth}\footnotesize
DFR pooled over decoders, module sizes and strengths at each error-correction level; the last
column is what moving from L to H buys. Redundancy helps most under occlusion and is not a
universal defence.
\end{minipage}
\end{table}
